\documentclass[12pt,a4paper]{report}

\usepackage[utf8]{inputenc}
\usepackage[T1]{fontenc}
\usepackage{graphicx}
\usepackage{hyperref}
\usepackage{setspace}
\usepackage{amsmath}
\usepackage{amssymb}
\usepackage[italian]{babel}
\usepackage[italian]{cleveref}

\hyphenation{sil-la-ba-zio-ne pa-ren-te-si}
\textwidth=450pt\oddsidemargin=0pt

\begin{document}

\title{Progettazione e Sviluppo di un Sistema di File Sharing}
\author{Alessandro Rebosio}

\maketitle

\pagenumbering{roman}
\setstretch{1.25}

\chapter*{Introduzione}
\label{chap:introduzione}
\addcontentsline{toc}{chapter}{Introduzione}
Negli ultimi anni, l'evoluzione delle tecnologie digitali ha
profondamente trasformato le modalità di gestione, archiviazione e
condivisione dei dati. In questo contesto, il Cloud Storage si è
affermato come una delle soluzioni più diffuse, consentendo agli
utenti di memorizzare informazioni su infrastrutture remote
accessibili tramite Internet. Servizi come Google Drive, Dropbox e
Microsoft OneDrive rappresentano esempi concreti di questo cambiamento,
offrendo strumenti efficienti, scalabili e semplici da utilizzare per
la gestione dei contenuti digitali.

Il cloud computing ha introdotto un cambiamento significativo
rispetto ai tradizionali sistemi di archiviazione locale. Invece di
dipendere esclusivamente da dispositivi fisici, come hard disk o
server aziendali, gli utenti possono accedere ai propri dati in
qualsiasi momento e da qualsiasi luogo, purché dispongano di una
connessione alla rete. Questo modello migliora la flessibilità
operativa e contribuisce a ridurre il rischio di perdita dei dati,
grazie a meccanismi di backup e ridondanza implementati dai fornitori
di servizi cloud.

Un aspetto particolarmente rilevante del Cloud Storage è il suo ruolo
nel favorire la collaborazione tra utenti. Le piattaforme moderne
permettono a più persone di lavorare simultaneamente sugli stessi
documenti, facilitando la condivisione in tempo reale e migliorando
l'efficienza dei processi lavorativi e accademici. La gestione dei
permessi e la sincronizzazione automatica contribuiscono a rendere
queste soluzioni strumenti fondamentali in contesti distribuiti.

Parallelamente, la persistenza dei dati rappresenta un elemento
cruciale nella progettazione di sistemi informatici moderni. Garantire
che le informazioni rimangano disponibili, integre e accessibili nel
tempo richiede l'adozione di architetture robuste e affidabili. In
questo contesto, assumono particolare importanza il modello di
storage adottato e le modalità con cui i diversi componenti del
sistema comunicano tra loro.

Alla luce di queste considerazioni, l'obiettivo del presente progetto
è la progettazione e realizzazione di un sistema di file sharing
scalabile, ispirato ai principi alla base delle moderne piattaforme
di Cloud Storage. Il sistema è stato sviluppato adottando
un'architettura a servizi containerizzati, al fine di garantire
modularità, isolamento delle componenti e facilità di deployment.

Il lavoro si propone inoltre di approfondire le principali
problematiche legate alla progettazione di sistemi distribuiti, con
particolare attenzione alla comunicazione tra container, alla gestione
dell'autenticazione e alla persistenza dei dati. Attraverso
l'implementazione del sistema e la valutazione delle sue
caratteristiche, si intende dimostrare come sia possibile realizzare
un'infrastruttura di file sharing che, pur in un contesto semplificato,
rifletta le logiche e i requisiti dei sistemi utilizzati in ambito
reale.

\tableofcontents

\cleardoublepage
\pagenumbering{arabic}
\setcounter{page}{1}

% CAPITOLO 1
\chapter{Obiettivi e Requisiti di Progetto}
\label{cap:obiettivi}

\section{Requisiti Funzionali}
% Descrivere le funzionalità offerte:
% - registrazione e autenticazione utenti
% - upload e download file
% - gestione dei file per utente
% Specificare attori, input e output.

\section{Requisiti Non Funzionali}
% Analizzare:
% - sicurezza (autenticazione, autorizzazione)
% - scalabilità (gestione carico e crescita utenti)
% - affidabilità (consistenza dati)
% - manutenibilità (modularità e separazione servizi)

\section{Vincoli Tecnologici}
% Descrivere:
% - uso di Docker per la containerizzazione
% - architettura a servizi separati
% - limite temporale (2 settimane)
% - tecnologie adottate (Node.js, React, PostgreSQL)

% CAPITOLO 2
\chapter{Architettura del Sistema}
\label{cap:architettura}

\section{Architettura a Container}
% Descrivere ogni container:
% - backend (Node.js): logica applicativa
% - database (PostgreSQL): persistenza dati
% - storage (filesystem o MinIO): gestione file
% - reverse proxy (Nginx): instradamento richieste
% - frontend (React): interfaccia utente
% Per ciascuno specificare:
% - ruolo
% - responsabilità
% - dipendenze

\section{Comunicazione tra Container}
% Descrivere:
% - rete Docker
% - comunicazione tramite hostname
% - porte interne vs esterne
% Spiegare i flussi principali:
% browser → Nginx → backend → DB / storage

\section{Infrastruttura e Networking}
% Approfondire:
% - ruolo di Nginx come reverse proxy
% - gestione delle richieste HTTP
% - separazione tra rete interna ed esterna

\section{Diagramma Architetturale}
% Inserire un diagramma che mostri:
% - tutti i container
% - connessioni tra di essi
% - flussi di dati
% Evidenziare chiaramente DB, backend, storage e proxy.

% CAPITOLO 3
\chapter{Containerizzazione e Deployment}
\label{cap:docker}

\section{Struttura dei Container}
% Descrivere:
% - immagini Docker utilizzate
% - configurazione dei servizi
% - dipendenze tra container

\section{Orchestrazione con Docker Compose}
% Spiegare:
% - come vengono avviati i servizi
% - definizione delle reti
% - collegamento tra container

\section{Gestione dei Volumi}
% Descrivere:
% - persistenza dati PostgreSQL
% - persistenza storage (filesystem / MinIO)
% - differenza tra volumi e storage effimero

\section{Vantaggi della Containerizzazione}
% Analizzare:
% - isolamento dei servizi
% - portabilità
% - facilità di deployment
% - confronto con architettura monolitica

% CAPITOLO 4
\chapter{Protocolli e Tecnologie Utilizzate}
\label{cap:protocolli}

\section{Comunicazione HTTP e REST}
% Descrivere:
% - uso di HTTP per le API
% - architettura REST
% - metodi principali (GET, POST)
% - upload file tramite multipart/form-data

\section{Comunicazione tra Servizi}
% Descrivere:
% - comunicazione backend → database
% - comunicazione backend → storage
% - uso di API S3 nel caso MinIO

\section{Persistenza dei Dati}
% Confrontare:
% - filesystem (struttura gerarchica)
% - object storage (bucket e oggetti)
% Evidenziare differenze in accesso e scalabilità.

% CAPITOLO 5
\chapter{Gestione Utenti e Autenticazione}
\label{cap:autenticazione}

\section{Modello di Autenticazione JWT}
% Descrivere:
% - generazione token
% - validazione
% - sessione stateless
% Inserire diagramma di sequenza del login.

\section{Sicurezza delle Credenziali}
% Spiegare:
% - hashing password (bcrypt)
% - autorizzazione accessi
% - protezione delle risorse

% CAPITOLO 6
\chapter{Implementazione v1: Storage su Filesystem}
\label{cap:v1_filesystem}

\section{Struttura delle Directory}
% Descrivere:
% - organizzazione file su disco
% - mapping DB ↔ filesystem

\section{Flusso Operativo}
% Descrivere il flusso completo:
% upload:
% browser → Nginx → backend → filesystem → DB
% download:
% backend → filesystem → utente
% Inserire diagramma di sequenza.

\section{Limiti della Soluzione}
% Analizzare:
% - scalabilità limitata
% - dipendenza dal nodo locale
% - gestione concorrente

% CAPITOLO 7
\chapter{Implementazione v2: Storage Remoto con MinIO}
\label{cap:v2_minio}

\section{Introduzione all'Object Storage}
% Descrivere:
% - concetto di bucket
% - oggetti
% - compatibilità con API S3

\section{Estensione dello Schema Dati}
% Presentare schema ER:
% - entità User
% - entità File
% - entità Bucket
% Relazioni:
% - User → File
% - File → Bucket
% Spiegare come cambia la gestione dei file.

\section{Integrazione con Storage Remoto}
% Descrivere:
% - uso SDK S3
% - operazioni upload/download/delete
% - gestione credenziali

\section{Flusso Operativo}
% Descrivere:
% upload:
% browser → Nginx → backend → MinIO → DB
% download:
% backend → MinIO → utente
% Evidenziare differenze rispetto alla v1.

% CAPITOLO 8
\chapter{Confronto e Valutazione delle Prestazioni}
\label{cap:confronto}

\section{Analisi Qualitativa}
% Confrontare filesystem vs object storage:
% - scalabilità
% - sicurezza
% - flessibilità
% - complessità

\section{Test di Prestazione}
% Descrivere:
% - metodologia
% - dimensioni file
% - tempi di risposta
% Confrontare v1 e v2.

% CAPITOLO 9
\chapter{Possibili Estensioni e Visione Futura}
\label{cap:estensioni}

\section{Gestione Permessi}
% Descrivere:
% - accessi ai file
% - condivisione tra utenti

\section{Scalabilità e Cloud}
% Descrivere:
% - migrazione verso cloud reale
% - sostituzione MinIO con servizi come Amazon S3
% - vantaggi e criticità

% CONCLUSIONI
\chapter*{Conclusioni}
\addcontentsline{toc}{chapter}{Conclusioni}
% Riassumere:
% - obiettivi raggiunti
% - differenze tra filesystem e object storage
% - vantaggi della containerizzazione
% - competenze acquisite (sistemi distribuiti, sicurezza, architettura)

\bibliographystyle{plain}
% \bibliography{report}
\end{document}

\end{document}
